% (C) 2026 Ralf Stephan, in collaboration with Claude Code.
% Released under CC0 1.0 Universal (public-domain dedication).
% See https://creativecommons.org/publicdomain/zero/1.0/

\documentclass[11pt]{article}
\usepackage[T1]{fontenc}
\usepackage[utf8]{inputenc}
\usepackage{lmodern}
\usepackage[margin=2.5cm]{geometry}
\usepackage{amsmath,amssymb}
\usepackage{booktabs,tabularx}
\usepackage{braket}
\usepackage[hidelinks]{hyperref}

% Lean names and paths: verbatim, breakable at . _ /
\DeclareUrlCommand\lean{\urlstyle{tt}}
\def\UrlBreaks{\do\.\do\_\do\/}
\newcolumntype{L}{>{\raggedright\arraybackslash}X}
\newcolumntype{B}{>{\raggedright\arraybackslash}p{0.42\textwidth}}
\emergencystretch=6em

\title{Lean formalization of an Assumptions of Physics technical brief:\\ log}
\author{}
\date{}

\begin{document}
\maketitle

\noindent
This log records which statements of the technical brief in this directory were checked in Lean,
where the formal version differs from the brief, and any problems found in the brief.
The brief is available at \url{https://assumptionsofphysics.org/resources/briefs}.

\section*{Setup}

\begin{itemize}
\item Lean \texttt{v4.36.0-rc1}; Mathlib at commit \texttt{2cabef97}; physlib (\texttt{master})
  and TauCeti (commit \texttt{eb4dd24}, on the same Mathlib) as further dependencies.
\item Build: \texttt{lake build} in the project root.
\item Every file is checked with \texttt{\#print axioms} on its main results. A result is
  recorded here only if it depends on no axioms beyond \lean{propext},
  \lean{Classical.choice} and \lean{Quot.sound}, and so has no \lean{sorry}.
\item Imports are kept to the Mathlib modules actually needed.
\item Material that belongs in Mathlib is kept in \lean{ForMathlib/}, in Mathlib's directory
  layout (library \lean{ForMathlib}), and imported by the brief files.
\end{itemize}


%======================================================================

\section{Quantum pure states are not enough}

\paragraph{Files.}\leavevmode
\begin{itemize}
\item File: \lean{briefs/PureStatesNotEnough/PureStatesNotEnough.lean}
\item Library: \lean{PureStatesNotEnough}
\item Imports: one new \lean{ForMathlib} file:\newline
  \lean{ForMathlib/Analysis/Convex/Combination.lean}: a map that preserves the convex combinations
  of two points of a convex set preserves all finite convex combinations
  (\lean{Convex.map_sum_of_map_add}), and a convex set contains them
  (\lean{Convex.sum_mem_of_semifield}), with scalars in any ordered semifield such as
  $\mathbb R_{\ge 0}$.\newline
  Existing files: \lean{ForMathlib/Analysis/InnerProductSpace/l2Space.lean}
  (\lean{Orthonormal.countable}), \lean{ForPhyslib/QuantumMechanics/DensityOperator.lean}
  (\lean{IsDensityOperator}, \lean{posTrace}), \lean{ForPhyslib/QuantumMechanics/Wigner/Density.lean}
  (\lean{IsDensity}, \lean{proj}) and
  \lean{ForPhyslib/ProbabilisticTheory/HilbertSpace/TraceClass/CanonicalExpansion.lean}
  (\lean{IsTraceClass.exists_hasSum_rankOne}).
\end{itemize}

\paragraph{Setting.} The ensembles of CST over a measurable space $X$ with Liouville measure
$\mu$ are the probability measures $\nu \ll \mu$ (\lean{classicalEnsembles}~$\mu$). Their mixtures
are $a\nu_1 + b\nu_2$ with $a, b \in \mathbb R_{\ge 0}$, $a + b = 1$.
\lean{HasManyDisjointSets}~$\mu$ says that for every $N$ there are $N$ disjoint measurable sets of
positive finite measure. The density operators of QE-DENS are \lean{IsDensity} (positive, trace
one) on a finite-dimensional space and \lean{IsDensityOperator} (positive, trace class, trace one)
on a separable complex Hilbert space $\mathcal H$ of any dimension. The pure states
(\lean{pureStates}) are the projections $|\psi\rangle\langle\psi|$ with $\|\psi\| = 1$, a copy
of $P(\mathcal H)$ inside the bounded operators, with the Borel $\sigma$-algebra of the operator
norm. An identification of two ensemble spaces is modelled as a map that preserves mixtures of two
ensembles.

\paragraph{Results.}
\begin{center}
\begin{tabularx}{\textwidth}{@{}BL@{}}
\toprule
Brief & Lean \\
\midrule
CST, the ensemble space & \lean{classicalEnsembles}; it is convex:
  \lean{Convex.classicalEnsembles} \\
(used for Theorem~3) a Darboux chart has $N$ disjoint sets of positive finite Liouville measure
  for every $N$ & \lean{hasManyDisjointSets_restrict}: any nonempty open set in a nontrivial
  finite-dimensional real normed space, for any measure that is positive on open sets and finite
  on compact sets \\
(used for Theorem~3) the CST ensemble space is infinite-dimensional &
  \lean{not_injOn_of_map_add}: no mixture-preserving map from it into a finite-dimensional real
  vector space is injective \\
Theorem~3, CST and QE-DENS are inconsistent & \lean{cst_qeDens_inconsistent}, for
  $\dim\mathcal H < \infty$ \\
Theorem~1, QP-RAY and CST are consistent iff $1 \le \dim V < \infty$ & not formalized;
  \lean{subsingleton_projectivization}: for $\dim V = 1$, $P(V)$ is a single point \\
Remark after Example~2, the Born rule is not $\delta_{\psi\phi}$ & \lean{born_ne_delta}: $e_0$ and
  $(e_0 + e_1)/\sqrt2$ are distinct rays with Born probability $\frac12$ \\
Theorem~5, the map $r : M^1(P(\mathcal H)) \to D(\mathcal H)$ & \lean{barycenter}
  $r(\mu) = \int \rho\,d\mu(\rho)$; \lean{isDensityOperator_barycenter},
  \lean{posTrace_barycenter}; \lean{toDensity} \\
Theorem~5, $r$ is affine & \lean{barycenter_add} \\
Theorem~5, $D(\mathcal H) \equiv M^1(P(\mathcal H))/{\sim}$ & \lean{toDensity_surjective},
  \lean{densityEquiv} \\
Remark after Theorem~5, a class records the different preparations of one ensemble &
  \lean{exists_barycenter_eq_and_ne}: $\frac12\delta_{e_0} + \frac12\delta_{e_1}$ and
  $\frac12\delta_{e_+} + \frac12\delta_{e_-}$ have the same barycenter but give $e_0$ different
  probabilities \\
\bottomrule
\end{tabularx}
\end{center}

\paragraph{Differences from the brief.}
\begin{itemize}
\item Theorem~3 compares the whole ensemble spaces, not those supported on a two-dimensional
  subspace $A$ (see ``Problems found''). It assumes $\dim\mathcal H < \infty$, which the brief
  derives from Theorem~1. Only the finite dimension of the real vector space of operators is
  used, not positivity or trace. ``The spaces correspond'' becomes a map from the CST ensembles to
  the density operators that is injective and preserves mixtures of two ensembles. The proof
  normalizes $\mu$ on $\dim + 2$ disjoint sets, takes an affine dependence among the images, and
  separates its positive and negative parts into two different mixtures with the same image.
\item The Liouville measure enters only through \lean{HasManyDisjointSets}, which holds on every
  Darboux chart. The symplectic structure itself is not used.
\item Theorem~5 is proved for every separable complex Hilbert space, not only in finite dimension.
  $r(\mu)$ is the Bochner integral of the pure states in the operator norm. It is positive because
  the positive operators form a closed convex cone, and it has trace one by Parseval and monotone
  convergence over a countable Hilbert basis. $r$ is onto by the canonical expansion
  $\rho = \sum_i\lambda_i|e_i\rangle\langle e_i|$, represented by $\sum_i\lambda_i\delta_{e_i}$.
  No Choquet theory is used. ``$\equiv$'' is the bijection between the quotient by
  $r(\mu_1) = r(\mu_2)$ and $D(\mathcal H)$.
\item Not formalized:
  \begin{itemize}
  \item Theorem~1;
  \item Example~2: $P(\mathbb C^2)$ is a sphere with Darboux coordinates $z$, $\theta_{xy}$;
    Hamiltonian flows; quantum evolution as rotations;
  \item Insight~4, which is informal;
  \item that the pure states are the extreme points of $D(\mathcal H)$, which the proof of
    Theorem~5 does not need.
  \end{itemize}
\end{itemize}

\paragraph{Problems found in the brief.} Theorems~3 and 5 hold, but their proofs have gaps.
Theorem~1 is not checked.
\begin{itemize}
\item Theorem~1 states $1 \le \dim V$, while its proof uses $2 \le \dim X$. For $\dim V = 1$ the
  state space $P(V)$ is a single point (\lean{subsingleton_projectivization}), a manifold of
  dimension $0$. So either $0$-dimensional symplectic manifolds count, and the proof's bound
  should be $0 \le \dim X$, or they do not, and the theorem should say $2 \le \dim V < \infty$.
  ``Since $X$ is a symplectic manifold, $\dim X < \infty$'' holds only for finite-dimensional
  symplectic manifolds, and the projective space of an infinite-dimensional Hilbert space is an
  infinite-dimensional one. The finiteness comes from CST itself: the Liouville measure $\omega^n$
  needs a finite $n$.
\item In the proof of Theorem~3, for $\dim\mathcal H \ge 3$ the sphere $P(A)$ has dimension
  $2 < 2(\dim\mathcal H - 1)$ and therefore Liouville measure zero. No CST ensemble has support in
  $P(A)$, so the classical side is empty, not infinite-dimensional. The argument works as written
  only for $\dim\mathcal H = 2$. Comparing the whole spaces repairs it.
\item The proof of Theorem~5 cites Choquet theory, which concerns compact convex sets, but
  $D(\mathcal H)$ is not compact when $\dim\mathcal H = \infty$. The proof also never shows that
  every density operator is represented, which $D(\mathcal H) \equiv M^1/{\sim}$ needs. Both gaps
  are closed by the canonical expansion. ``Every continuous linear functional
  $f : D(\mathcal H) \to \mathbb R$'' should read continuous affine function, since
  $D(\mathcal H)$ is not a vector space.
\item Example~2 (not checked in Lean): area-preserving diffeomorphisms of the sphere include the
  orientation-reversing ones, such as $z \mapsto -z$, which are not Hamiltonian. So Hamiltonian
  evolution preserves the area form, not only the area.
\item Typos: ``non necessarily'' should be ``not necessarily''; ``$r$ defined about'' should be
  ``$r$ defined above''.
\end{itemize}

\end{document}
